\section{Framework, episodes and method}\label{sec:framework}

\subsection{The estimand}

For every episode the object is the same conditional quantity: the
response of the statutory tax--benefit system to a specified shock.
A shock is (i) a vector of consumer-price changes by consumption
category, or (ii) a labour-income loss with a sectoral location and an
adjustment margin. The household stage maps it onto households and
computes the change in disposable income, its distribution and the
Exchequer's share, under the rules of a stated year
(Section~\ref{sec:secondstage-method}).

The headline metric is the \emph{cushioning rate}. For an earnings
shock it is the income-stabilisation coefficient of \citet{dolls2012},
\begin{equation*}
c^{Y} \;=\; 1 - \frac{\Delta Y^{\text{disp}}}{\Delta Y^{\text{mkt}}},
\end{equation*}
one minus the change in disposable income per pound of market-income
loss. For a price shock disposable income does not move, so the
analogue is
\begin{equation*}
c^{P} \;=\; \frac{T}{\Delta E}, \qquad
\Delta E \;=\; \sum_{k} p_{k} q_{k}\,\hat{p}_{k},
\end{equation*}
where $\Delta E$ is the cost of the price changes $\hat{p}$ at fixed
consumption and $T$ is the support the episode triggered, automatic
plus discretionary. The two are analogues, not the same parameter:
$c^{Y}$ is the share of a market-income loss absorbed by lower taxes
and higher benefits; $c^{P}$ is the ratio of legislated support to the
fixed-basket cost of the price change, and it can exceed one because
lump-sum payments are not tied to the loss. Episodes are compared in
such ratios, and in burdens as shares of spending, because their sizes
span more than two orders of magnitude and pound totals would reduce
any comparison to a statement about the energy episode.

\subsection{Two channels and a policy-year comparison}

The expenditure channel prices a category-level price vector against
household expenditure shares, first on published
expenditure-by-decile tables and then on imputed household
consumption. The earnings channel, the labour-income channel of
\citet{porto2006}, imposes sectoral labour-income changes on
employees by industry and runs them through the tax--benefit system,
as wage cuts or as displacement; the imposed-transition step follows
\citet{brewertasseva2021} and \citet{navicke2014}. The channels are
reported separately and not summed: a price shock erodes real income
that the benefit system indexes with a lag, while an earnings shock
changes tax liabilities and means-tested entitlement at once. The two
responses run through different statutory mechanisms, and their
comparison is the point of setting the episodes side by side.

Benefit uprating follows lagged inflation, so a fast price shock
outruns indexation for a year, and discretionary interventions rewrite
the schedule mid-shock. For the energy episode the paper therefore
separates three quantities: the cushioning delivered by the rules as
they stood; the benefit income the same rules would have paid had the
April 2022 uprating matched that year's inflation; and the
discretionary instruments on top. Price-shock incidence has been run
through tax--benefit models before, and indexation counterfactuals are
standard in that literature (Section~\ref{sec:literature}); what is
added here is the separation of the statutory timing lag from the
discretionary instruments within one episode, on the same households,
and the split of those instruments into energy-specific and general
support.

\subsection{Scope and interpretation}\label{sec:framework-ident}\label{sec:framework-decile}

Three points govern how the results should be read. First, nothing
here estimates a shock. Whether the Trade and Cooperation Agreement
caused the food-price rise is identified, or not, by the
quasi-experimental design of \citet{bakker2023}; the UK--India figures
are the government's model output; the tariff earnings loss is
assumed; and the energy shock is observed in regulated cap levels,
with the share that is trade a matter of Ofgem's cost accounting
(Appendix~\ref{app:wholesale}). Given a shock, what the system pays a
household is a calculation on statutory rules, though it inherits
measurement error from imputed consumption and simulated entitlement.
The results are conditional incidence: given a shock of stated size,
how its cost fell across households after taxes, benefits and
legislated support.

Second, the estimates are first-round and static. There is no
behavioural response, no wage or employment feedback from the price
shocks and no substitution in consumption, so the expenditure-channel
numbers are first-order approximations that overstate the welfare loss
where households substitute. Each shock is carried at its central
value. Only the consumption-basket channel of the price shocks is
measured; \citet{ferreira2026} find the net-nominal-position and
labour-income channels of inflation larger.

Third, decile means describe whether a shock was regressive, not who
was hurt. \citet{borusyakjaravel2024} and \citet{levell2025} show
that most of the dispersion in exposure lies within income deciles,
which Section~\ref{sec:results-withindec} confirms for energy on
household records.
